\begin{table}[!htbp]
\caption{整数计数预测系统中新增城市数据源的实际决策框架。}
\label{tab:practical_decision_framework}
\centering
\footnotesize
\begin{tabularx}{\textwidth}{p{2.8cm}XXp{3.2cm}}
\toprule
\textbf{决策问题} & \textbf{所需证据} &
\textbf{操作响应} & \textbf{边界}\\
\midrule
数据源是否产生估计预测机会？ &
冻结目标总体中的 DELB 与 CMI 点估计 &
若估计底线实质下降，则保留该数据源进入校准评价 &
点估计不能建立数据源特异信息或实际收益\\
表观增益能否与有限样本结构区分？ &
支持度诊断、已知真值校准与设计匹配随机化 &
仅在校准和功效支持分离时，才对所检验设计进行有限归因 &
低功效下不拒绝没有确定性；随机化不是因果识别\\
现有模型能否在未见数据上利用该数据源？ &
目标、样本、调参和损失均匹配的留出模型 &
只部署表现出稳定留出收益的城市--模型组合，并监测 Predictability Gap &
单一模型改善不能推广到其他城市或模型族\\
价值是否延伸到服务覆盖区外？ &
训练期冻结的服务覆盖总体与完整目标总体 &
同时报告两个总体，并把覆盖范围视为数据产品的一部分 &
服务区估计不能外推到未覆盖地点\\
\bottomrule
\end{tabularx}
\end{table}
